\scanpage{261}

\mspage{133}

donc ils sont congrus dans $T_G$ modulo $\unsure{\mathfrak{z}|\mathfrak{Z}}$. Pour \deleted{voir} qu'ils soient égaux, il \deleted{suffit de} \inserted{\unsure{f. et s.|suffit}} que leurs images dans $G/H\simeq\widetilde Z$ soient les $=$s, puisque l'on a $\mathfrak{z}\cap H=\{1\}$, i.e. que leurs images par $\widetilde\delta$ soient les ms, ce qui s'écrit sous la forme

\[c'(c(\widetilde\lambda))=\widetilde\lambda^2\qquad\forall\widetilde\lambda\in\widetilde Z\tag*{(133)}\]

(à rapprocher de (127), qui implique que les deux mbes de (\unsure{133|138}) sont égaux mod $\mathfrak{z}$).

En résumé

\[\begin{aligned}
&\forall\widetilde\lambda\in\widetilde Z,\ \text{on a}\\
&\hspace{3em}\underset{\substack{\rotatebox{90}{$\in$}\\Z\subset T_G}}{c(\widetilde\lambda)}
\equiv u_0(\widetilde\lambda)u_1(\widetilde\lambda)\quad\bmod\mathfrak{z}.
\end{aligned}\tag*{(134)}\]

Pour que cette congruence soit une égalité, il f. et s. que l'on ait (133), (qui implique (131)).

Quand de plus \deleted{\unsure{si}} $c:\widetilde Z\unsure{-|\longrightarrow}Z$ est surjectif, la formule (133) redonne une façon d'expliciter le plongement de $Z$ dans $T_G$, en termes de la seule donnée du diagramme (126).

En résumé, pour reconstituer tout le diagramme (64) à isom près -- en même temps que (113), -- à partir des données (\unsure{89|88})

% archive sheet number: 260
